%% toc: P15 166 166.1
% Proofs for the discrete entropy functional and Shannon's characterization theorem.
\section{Entropy of discrete distributions: proofs}
\label{sec:annex-p15-w02b-p15-040-information}

Use the definitions of \cref{sec:p15-w02b-p15-040-discrete}. Logarithms
are natural. Two elementary facts are used repeatedly.

\emph{(L) For every \(u>0\), \(\log u\leq u-1\), with equality only at
\(u=1\).} Put \(g(u)=u-1-\log u\). Then \(g'(u)=1-1/u\) is negative
on \((0,1)\) and positive on \((1,\infty)\), and \(g(1)=0\). By the
mean value theorem \(g\) is strictly decreasing on \((0,1]\) and
strictly increasing on \([1,\infty)\), so \(g(u)>0\) for \(u\ne1\).

\emph{(M) Shannon's function \(A\) is nondecreasing under (S2), and
under its weakened form.} Either form gives \(A(n)\leq A(n+1)\) for all
\(n\), and chaining these inequalities gives \(A(a)\leq A(b)\) whenever
\(1\leq a\leq b\).

\begin{annexitem}{annex:p15.w02b-p15-040.domain}{thm:p15-w02b-p15-040-domain}{Well-definedness, units and extreme values}
\begin{proof}
\emph{(1)} Let \(\varphi(t)=-t\log t=t\log(1/t)\) for \(0<t\leq1\).
Since \(1/t\geq1\), \(\varphi(t)\geq0\). Its derivative
\(\varphi'(t)=-\log t-1\) is positive on \((0,1/e)\) and negative on
\((1/e,1]\), so \(\varphi\) attains its maximum \(\varphi(1/e)=1/e\).
Thus every term of the defining sum lies in \([0,1/e]\). The sum of a
family of nonnegative extended reals indexed by a countable set is the
supremum of its sums over finite subsets. That supremum exists in
\([0,+\infty]\) and equals the limit of the partial sums along any
enumeration, since every finite subset is contained in some initial
segment and every partial sum is a finite-subset sum. Adding or deleting
indices with \(p_i=0\) adds or deletes terms equal to zero under the
definition, which sums only over \(p_i>0\). If \(\sigma:I'\to I\) is a
bijection, \(F\mapsto\sigma^{-1}(F)\) is a bijection between the finite
subsets of \(I\) and of \(I'\) that preserves each finite sum, so the
suprema agree.

\emph{(2)} On \((0,1]\), \(\varphi\) is continuous. For
\(0<t\leq1\), (L) with \(u=t^{-1/2}\) gives
\(\log(1/t)=2\log t^{-1/2}\leq2(t^{-1/2}-1)<2t^{-1/2}\), hence
\[
 0\leq\varphi(t)<2\sqrt t\longrightarrow0\qquad(t\downarrow0).
\]
So \(\varphi\) is continuous at \(0\) with \(\varphi(0)=0\), which is
the convention \(0\log0=0\). On \(\Delta_n\),
\(H(p)=K\sum_{i=1}^n\varphi(p_i)\) is a finite sum of continuous
functions of the coordinates, each at most \(1/e\); it is finite and
continuous.

\emph{(3)} For \(u>0\), \(\log_bu=\log u/\log b\), with
\(\log b>0\) because \(b>1\). Summing over \(p_i>0\),
\(H_b(p)=\sum p_i\log(1/p_i)/\log b=H_e(p)/\log b\). The functional
with constant \(K\) is \(KH_e\); it equals \(H_e/\log b=H_b\) exactly
when \(\log b=1/K\), that is \(b=e^{1/K}\), and \(e^{1/K}>1\) because
\(K>0\). Conversely \(H_b\) is \(KH_e\) with \(K=1/\log b>0\). Since
\(\log2\approx0.693147\), \(1/\log2\approx1.442695\).

\emph{(4)} By (1), \(H(p)=0\) if and only if every term
\(p_i\log(1/p_i)\) with \(p_i>0\) is zero. For \(0<t<1\),
\(\log(1/t)>0\), so a term vanishes only when \(p_i=1\). Hence
\(H(p)=0\) if and only if every positive \(p_i\) equals \(1\); since the
\(p_i\) sum to \(1\), this means exactly one \(p_i\) equals \(1\) and
the rest are zero. If \(p\) is not of this form, some term is positive
and \(H(p)>0\), possibly \(+\infty\) in the countable case.

\emph{(5)} Apply \cref{thm:p15-rmdl-gibbs} to the masses \(p\) and
\(q_i=1/n\) on \(\{1,\ldots,n\}\), both of total \(1\). No \(q_i\) is
zero, and
\[
 0\leq D(p,q)=\sum_{i:\,p_i>0}p_i\log(np_i)
   =\log n-\sum_{i:\,p_i>0}p_i\log\frac1{p_i}=\log n-H_e(p),
\]
with equality if and only if \(p=q\). Multiplying by \(K>0\) gives
\(H(p)\leq K\log n\), with equality exactly at the uniform vector.

\emph{(6)} Let \(I\) be countably infinite. For each \(n\), the vector
equal to \(1/n\) on \(n\) chosen indices and \(0\) elsewhere lies in
\(\Delta(I)\) and, by (1) and (5), has entropy \(K\log n\). These values
are unbounded.

For the example, put \(g(x)=1/(x(\log x)^2)\) and
\(f(x)=1/(x\log x)\) for \(x>1\). Both are positive, and both are
strictly decreasing because \(x\), \(\log x\) and \((\log x)^2\) are
positive and strictly increasing on \((1,\infty)\). For an integer
\(k\geq3\), \(g(k)\leq\int_{k-1}^kg(x)\,dx\), so for \(N\geq3\)
\[
 \sum_{k=3}^Ng(k)\leq\int_2^N\frac{dx}{x(\log x)^2}
   =\frac1{\log2}-\frac1{\log N}<\frac1{\log2},
\]
since \(-1/\log x\) has derivative \(g(x)\). Hence
\(S\defeq\sum_{k\geq2}g(k)\leq g(2)+1/\log2<\infty\). Also
\(S\geq g(2)=1/(2(\log2)^2)>1\), because \(\log2<0.7\) gives
\((\log2)^2<0.49<\tfrac12\). Thus \(c=1/S\) satisfies \(0<c<1\), each
\(p_k=cg(k)\) is positive, and \(\sum_{k\geq2}p_k=cS=1\).

For \(k\geq2\),
\[
 \log\frac1{p_k}=\log k+2\log\log k+\log\frac1c .
\]
For \(k\geq3\), \(\log k\geq\log3>1\) because \(3>e\), so
\(\log\log k>0\); and \(\log(1/c)>0\) because \(c<1\). Therefore
\(p_k\log(1/p_k)\geq p_k\log k=cf(k)\) for \(k\geq3\). For \(k\geq3\),
\(f(k)\geq\int_k^{k+1}f(x)\,dx\), and \(\log\log x\) has derivative
\(f(x)\), so for \(N\geq3\)
\[
 \sum_{k=3}^Nf(k)\geq\int_3^{N+1}\frac{dx}{x\log x}
   =\log\log(N+1)-\log\log3\longrightarrow+\infty .
\]
Since all terms are nonnegative,
\(H(p)\geq Kc\sum_{k\geq3}f(k)=+\infty\).
\end{proof}
\end{annexitem}

\begin{annexitem}{annex:p15.w02b-p15-040.characterization}{thm:p15-w02b-p15-040-characterization}{Shannon's Theorem~2 and Appendix~2}
The proof follows Shannon's Appendix~2
\cite[pp.~28--29]{Shannon1948MathematicalTheory}. Steps~1 to~4 use only
(S1), the nondecreasing property (M) and the uniform instances (S3u).
\begin{proof}
\emph{Step 0: \(A(1)=0\).} Apply (S3u) with \(n=1\), \(p=(1)\),
\(m_1=m\) and \(q^{(1)}\) uniform on \(m\) outcomes. The composite vector
is \(q^{(1)}\) itself, so \(A(m)=H_1(1)+1\cdot A(m)=A(1)+A(m)\). Since
\(A(m)\) is a real number, \(A(1)=0\).

\emph{Step 1: \(A(s^m)=mA(s)\) for integers \(s\geq1\), \(m\geq0\).}
For \(m=0\) this is Step~0. Suppose it holds for \(m\). Apply (S3u) with
\(n=s^m\), \(p\) uniform on \(s^m\) outcomes, and every \(m_i=s\) with
\(q^{(i)}\) uniform. Each of the \(s^{m+1}\) composite entries equals
\(s^{-m}s^{-1}\), so the composite vector is uniform on \(s^{m+1}\)
outcomes and
\[
 A(s^{m+1})=A(s^m)+\sum_{i=1}^{s^m}s^{-m}A(s)=mA(s)+A(s).
\]
This is Shannon's decomposition of a choice from \(s^m\) equally likely
possibilities into \(m\) successive choices from \(s\).

\emph{Step 2: \(A(t)=K\log t\).} First assume \(A(2)>0\). Then
\(A(s)\geq A(2)>0\) for every \(s\geq2\) by (M). Fix integers
\(s,t\geq2\) and \(n\geq1\). Since \(s^0=1\leq t^n\) and \(s^m\to\infty\),
there is a largest integer \(m\geq0\) with \(s^m\leq t^n\), and then
\[
 s^m\leq t^n<s^{m+1}.
\]
Taking logarithms and dividing by \(n\log s>0\),
\[
 \frac mn\leq\frac{\log t}{\log s}<\frac mn+\frac1n .
\]
By (M) and Step~1,
\(mA(s)=A(s^m)\leq A(t^n)=nA(t)\leq A(s^{m+1})=(m+1)A(s)\), and
dividing by \(nA(s)>0\),
\[
 \frac mn\leq\frac{A(t)}{A(s)}\leq\frac mn+\frac1n .
\]
Both ratios lie in an interval of length \(1/n\), so
\(\bigl|A(t)/A(s)-\log t/\log s\bigr|\leq1/n\). As \(n\) is arbitrary,
\(A(t)/\log t=A(s)/\log s\) for all \(s,t\geq2\). With
\(K\defeq A(2)/\log2>0\), \(A(t)=K\log t\) for \(t\geq2\), and
\(A(1)=0=K\log1\).

Under the strict requirement (S2), \(A(2)>A(1)=0\), so this case
applies and \(K>0\). Under the weakened requirement, it remains to treat
\(A(2)=0\). Then \(A(2^m)=0\) for every \(m\) by Step~1. For any
\(t\geq1\), \(1\leq t<2^t\), so by (M)
\(0=A(1)\leq A(t)\leq A(2^t)=0\). Thus \(A(t)=0=K\log t\) with \(K=0\).

\emph{Step 3: the formula at rational vectors with positive entries.}
Let \(p\in\Delta_n\) with \(p_i=n_i/N\), where the \(n_i\geq1\) are
integers and \(N=\sum_in_i\). Apply (S3u) with \(m_i=n_i\) and
\(q^{(i)}\) uniform on \(n_i\) outcomes. Each composite entry is
\(p_i/n_i=1/N\), and there are \(N\) of them, so
\[
 K\log N=A(N)=H_n(p)+\sum_{i=1}^np_iA(n_i)=H_n(p)+K\sum_{i=1}^np_i\log n_i .
\]
Since \(\sum_ip_i=1\),
\[
 H_n(p)=K\Bigl(\sum_ip_i\log N-\sum_ip_i\log n_i\Bigr)
   =-K\sum_ip_i\log\frac{n_i}N=-K\sum_ip_i\log p_i .
\]
This is Shannon's computation of the total choice from \(\sum n_i\)
possibilities by two methods.

\emph{Step 4: all vectors, by continuity.} Let \(p\in\Delta_n\). For
\(N\geq1\) put \(n_i^{(N)}=\lfloor Np_i\rfloor+1\), an integer at least
\(1\), and \(r^{(N)}_i=n_i^{(N)}/\sum_jn_j^{(N)}\). Then
\(Np_i<n_i^{(N)}\leq Np_i+1\), so \(n_i^{(N)}/N\to p_i\) and
\(\sum_jn_j^{(N)}/N\to\sum_jp_j=1\); hence \(r^{(N)}\to p\) in
\(\mathbb R^n\). Each \(r^{(N)}\) is a rational vector with positive
entries, so by Step~3
\(H_n(r^{(N)})=-K\sum_ir^{(N)}_i\log r^{(N)}_i\). The left side
tends to \(H_n(p)\) by (S1); the right side tends to
\(-K\sum_ip_i\log p_i\) by item~2 of
\cref{thm:p15-w02b-p15-040-domain}. This proves items~1 and~3; in
item~1, \(K=A(2)/\log2\) by Step~2. Shannon approximates
incommensurable probabilities by rationals and appeals to continuity in
the same way.

\emph{Item 2.} Let \(K>0\) and \(H_n(p)=K\sum_i\varphi(p_i)\), with
\(\varphi\) as in the proof of \cref{thm:p15-w02b-p15-040-domain}.
(S1) is item~2 of that theorem. For (S2),
\(A(n)=K\sum_{i=1}^n\tfrac1n\log n=K\log n\), which is strictly
increasing in \(n\). For (S3), first note that for \(a,b\in[0,1]\),
\[
 \varphi(ab)=a\varphi(b)+b\varphi(a).
\]
If \(a,b>0\) this is \(-ab\log(ab)=-ab\log b-ab\log a\); if \(a=0\) or
\(b=0\) both sides are zero because \(\varphi(0)=0\). The composite
vector in (S3) has nonnegative entries summing to
\(\sum_ip_i\sum_jq^{(i)}_j=1\), so it lies in \(\Delta_M\), and
\[
 \sum_{i=1}^n\sum_{j=1}^{m_i}\varphi\bigl(p_iq^{(i)}_j\bigr)
 =\sum_{i=1}^np_i\sum_{j=1}^{m_i}\varphi\bigl(q^{(i)}_j\bigr)
  +\sum_{i=1}^n\varphi(p_i)\sum_{j=1}^{m_i}q^{(i)}_j
\]
\[
 =\sum_{i=1}^np_i\frac{H_{m_i}(q^{(i)})}K+\frac{H_n(p)}K .
\]
Multiplying by \(K\) gives (S3).
\end{proof}
\end{annexitem}
